\subsubsection{Synaptic Intelligence (SI)}

\gls{si} is a regularisation-based \gls{cl} method that estimates parameter
importance online from the optimisation trajectory. Rather than computing a
post-hoc curvature estimate, \gls{si} accumulates a path integral that measures
how much each parameter contributed to reducing loss during a task.

\paragraph{Path-Integral Importance}

During training on task $t$, \gls{si} tracks a contribution accumulator
$W_i^{(t)}$ for each parameter $\theta_i$:
%
\begin{equation}\label{eqn:si_path}
    W_i^{(t)}
    \leftarrow
    W_i^{(t)} - g_i \,\Delta\theta_i
\end{equation}
%
where $g_i$ is the gradient and $\Delta\theta_i$ is the parameter change at the
step. At the task boundary, this path integral is normalised by the total
parameter displacement to obtain importance:
%
\begin{equation}\label{eqn:si_omega}
    \Omega_i
    \leftarrow
    \Omega_i
    +
    \frac{W_i^{(t)}}{(\theta_i-\theta_i^*)^2+\epsilon}
\end{equation}
%
where $\theta_i^*$ is the parameter value stored at the previous
consolidation point and $\epsilon$ stabilises the denominator.

\paragraph{\gls{cl} Approach}

When learning subsequent tasks, \gls{si} augments current-task cross-entropy
with a quadratic consolidation penalty:
%
\begin{equation}\label{eqn:si_loss}
    \mathcal{L}_\text{SI}
    =
    \mathcal{L}_\text{CE}
    +
    c
    \sum_i \Omega_i(\theta_i-\theta_i^*)^2
\end{equation}
%
where $c$ controls regularisation strength. Parameters with high accumulated
importance are constrained to remain near prior consolidated values, while
low-importance parameters remain adaptable for new tasks.

\paragraph{Inference}

\gls{si} uses a standard linear classifier at inference time. In \gls{til}, task
identity selects the corresponding task logit range. In \gls{cil}, prediction
is made from the full logit vector across all seen classes.

